\section{Ausblick}
\label{sec:Ausblick}

Die Ergebnisse eröffnen Weiterentwicklungspotenziale. Technisch ermöglicht die Evaluation größerer lokaler Modelle (13B, 70B) weitere Qualitätssteigerungen. Dies setzt erweiterte Hosting-Infrastrukturen voraus. Die Integration komplexer Strukturen wie itemAssociationRules oder variable Prozessformeln ist bereits validiert. Diese kann systematisch ausgebaut werden. Automatische Dateisystem-Scanner für ERP/MES-Datenextraktion würden den Self-Service-Ansatz weiter stärken. Die Erweiterung auf zusätzliche Entitätsstrukturen sollte als nächstes angestrebt werden. Dies ermöglicht komplexere Konfigurationen. Der Self-Service-Ansatz mit Master Creator ermöglicht Skalierung auf größere Kundenzahlen. Eine proportionale Teamvergrößerung ist nicht erforderlich. Die strategische Flexibilität optimiert Time-to-Market. Drei Nutzungsformen stehen zur Verfügung: reine Master Creator-Nutzung, iterative Entwicklung oder nachgelagerte Importer-Integration. Diese passen sich unterschiedlichen Kundenkapazitäten und Datenlagen an.

Die konzeptuelle Trennung zwischen domänenunabhängiger Methodik und GRACE-spezifischen Inhalten legt Transferierbarkeit nahe. Die Drei-Artefakte-Architektur ist auf andere Enterprise-Systeme übertragbar. Fragebogen, Entity Mapping und SOPs können inhaltlich angepasst werden. Anwendungsfelder umfassen SAP-Customizing, MES-Setup oder CRM-Konfiguration. Eine weiterführende Vision erschließt sich durch Verknüpfung dreier Automatisierungsebenen. Erste Ebene: Automatisierte Datenerhebung durch COSIMA (Datenraumadapter, aktuell in Entwicklung bei seitcom GmbH). Kundensysteme werden auf Dateisystemebene nach APS-relevanten Informationen durchsucht und strukturiert. Zweite Ebene: Automatisierte Konfiguration durch Master Creator. Dieser generiert GRACE-Entitäten konversationell aus strukturierten Daten. Dritte Ebene: Automatisierte APS-Planung durch GRACE. Das System führt KI-gestützte Produktionsplanung auf Basis der generierten Konfigurationen durch. Diese End-to-End-Automatisierung würde den gesamten APS-Lebenszyklus abdecken. Der Zyklus reicht von der initialen Datenerhebung über die Konfigurationsgenerierung bis zur operativen Planberechnung. Die Abhängigkeit von OR-Consultants würde auf Qualitätssicherung und Sonderfallerkennung reduziert.
